%\vspace{-0.5em}

% \paragraph{Time series forecasting with switching dynamical systems}  State-space models (Fraccaro et al., 2017; Rangapuram
% et al., 2018; de Bézenac et al., 2020; Tang & Matteson, 2021; Klushyn et al., 2021) and recurrent neural
% networks (Salinas et al., 2019; 2020; Rasul et al., 2021b;a) are among the most popular choices to model
% temporal dependencies in time series. 

% However, for these models to generalize and adapt to new changes in the future, it is
% critical to have appropriate inductive biases built into the model. A common approach is to allow state-space models to switch between a discrete set of dynamics, which can be learned from training data.

\vspace{-0.5em}
\paragraph{Probabilistic Forecasting for Time-Series with Change Points.} %Deep learning has enabled substantial advances in time series forecasting. in order to represent the inherent stochasticity of real-world time series data. 
%Probabilistic forecasting methods model distributions over possible outcomes \citep{benidis2022deep}, and have been used widely to build predicting intervals.
%Bayesian Neural Networks (BNN) for uncertainty quantification (UQ) of spatiotemporal forecasts. In practice, Bayesian UQ can be computationally expensive and difficult to optimize, especially for larger networks \citep{lakshminarayanan2017simple, zadrozny2001obtaining}. Therefore, UQ for deep neural network time series forecasts often adopts approximate Bayesian inference such as MC-dropout \citep{gal2016theoretically, gal2017concrete}. Neural process models use neural network to parameterize a stochastic process \citep{garnelo2018conditional,garnelo2018neural, sun2024data}. 
%Method such as (approximate) Bayesian Neural Networks  \citep{wang2019deep, wu2021quantifying, gal2017concrete} and neural process models \citep{garnelo2018conditional, sun2024data} use neural networks to parametrize traditional probabilistic frameworks for time-series forecasting. Other deep-learning models such as DeepAR \citep{salinas2020deepar}, PatchTST \citep{nie2023patchtst}, and Temporal Fusion Transformer \citep{lim2021temporal} directly generate probabilistic forecasts by training the model with pinball loss. 

Probabilistic forecasting has become central to modern time-series analysis by modeling a distribution over future outcomes rather than a single-point forecast \citep{benidis2022deep}. Classic approaches often adopt Bayesian or approximate Bayesian strategies for uncertainty quantification: for instance, Bayesian Neural Networks (BNNs) \citep{wang2019deep, wu2021quantifying, gal2017concrete} or neural process models \citep{garnelo2018conditional, sun2024data} use neural networks to parameterize the underlying stochastic process. Other methods, such as DeepAR \citep{salinas2020deepar}, PatchTST \citep{nie2023patchtst}, or Temporal Fusion Transformer \citep{lim2021temporal}, generate probabilistic forecasts directly by minimizing a pinball loss (i.e., quantile loss), thereby learning predictive distributions.

%\textit{Change points} refer to abrupt distribution shifts in observational time series data \citep{montanez2015inertial}. Detecting and segmenting change points has been studied extensively (see \cite{aminikhanghahi2017survey} for a recent survey), %\citep{adams2007bayesian, agudelo2020bayesian, chiappa2010movement}, 
%whereas modeling time series with change points is typically approached with SSM, and in particular SDS. Recent works in SDS 
%mostly deploy variational inference \citep{johnson2016composing, ansari2021deep, liu2023graph}, using neural networks to parametrize the transition or emission distributions. \cite{sun2024online} studies UQ for change points but in the anomaly detection setting. % to improve flexibility and accelerate inference.

A key challenge arises when \textit{change points} occur - i.e., abrupt distribution shifts in the time series \citep{montanez2015inertial}. Change point detection and segmentation have been studied extensively (beyond our scope, see \citep{aminikhanghahi2017survey} for a recent survey). Modeling time series with change points typically features state-space models (SSM), and in particular switching dynamical systems (SDS) \citep{johnson2016composing, ansari2021deep, liu2023graph}. These SDS approaches use variational inference to fit neural parameterizations of transitions and emission distributions. However, existing probabilistic forecasters often lack strict calibration guarantees and can be overconfident in practice \citep{lakshminarayanan2017simple}, and many require specialized architecture tuning for different tasks, limiting their applicability to real-world UQ challenges.


% Compared to ours and other CP approaches, the drawbacks of the aforementioned methods are: (1) they are not model-free, i.e. special designs of the predictive model and hyper-parameter tuning are required, and (2) they have no theoretical guarantee on calibration and often are overconfident in practice \citep{lakshminarayanan2017simple}.

%SDS models, most often represented as a graphical model, have seen a surge in using neural networks to parameterize transition and emission distributions \citep{johnson2016composing, ansari2021deep, liu2023graph}, to improve flexibility and fast inference. Bayesian and variational inference methods are used extensively in online change point detection \citep{adams2007bayesian, agudelo2020bayesian} and segmentation \citep{chiappa2010movement} as well. 

% Our work focuses on only discrete states, which come with the benefit of interpretability and scalability. Discrete SSMs have been used to address safety-critical distribution shifts, such as reasoning about occlusion in traffic scenes \cite{wei2024navigation}. 

%The challenge of these Bayesian, RNN-, or Transformer-based probabilistic models is that the forecasted probabilities are often uncalibrated or overconfident \citep{lakshminarayanan2017simple}, hence unsuitable for many practical tasks that require guarantees on the uncertainty estimates associated with the forecasts. 

%Frequentist UQ methods emphasize robustness against variations in the data. These approaches either rely on resampling the data or learning an interval bound to encompass the dataset. For time series forecasting UQ, approaches include ensemble methods such as bootstrap \citep{efron2016, bjrnn} and jackknife methods  \citep{jackknife+after, bjrnn}; interval prediction methods include interval regression through proper scoring rules \citep{kivaranovic2020adaptive, wu2021quantifying}, and quantile regression \citep{takeuchi2006nonparametric}, with many recent advances for time series UQ  \citep{tagasovska2019single, gasthaus2019probabilistic,park2022learning, kan2022multivariate}. Many of the frequentist methods produce asymptotically valid confidence regions and can be categorized as distribution-free UQ techniques as they are (1) agnostic to the underlying model and (2) agnostic to data distribution.

\vspace{-0.5em}
\paragraph{Conformal Prediction for time-series.} 
%Conformal prediction (CP), originally introduced in \cite{vovk2005algorithmic}, has become a popular UQ method because of its simplicity, generality,  coverage guarantee, and low computational cost; we refer readers to \cite{angelopoulos2021gentle} for a comprehensive introduction to CP. For time series prediction, \cite{Schaar2021conformaltime, sun2023copula, zhou2024heterogeneous} aims for joint coverage and formulates fixed-length time series problems as exchangeable multi-target forecasts. Enbpi \citep{xu2021conformal} and SPCI \citep{xu2023sequential} assume no change in system dynamics, but in the generating process of the noise or residual that causes uncertainty.  \cite{tibshirani2019conformal} adapts to \textit{known} covariance shift by weighting data using propensity scores. Although theoretically optimal, the algorithm does not apply to real-world problems where propensity scores are difficult to calculate. Adaptive conformal inference (ACI) \citep{gibbs2021adaptive} and its variants \citep{zaffran2022adaptive, gibbs2024conformal, angelopoulos2024pid} consider the online adversarial setting, where they guarantee target marginal coverage without any assumptions on the data generating process. 

Conformal prediction (CP) \citep{vovk2005algorithmic} has emerged as a leading framework for UQ due to its algorithmic simplicity, broad applicability, and finite-sample coverage guarantees; see \cite{angelopoulos2021gentle} for a comprehensive introduction. For time-series forecasting, some works focus on joint coverage over fixed-length horizons \citep{Schaar2021conformaltime, sun2023copula, zhou2024heterogeneous}; others assume stable underlying dynamics and adapt only to shifts in residual distributions \citep{xu2021conformal} . \cite{tibshirani2019conformal} uses known covariate shifts (by propensity weighting) to achieve theoretical coverage, but can fail in settings where the shift mechanism is unknown, which is typical for time series forecasting. \cite{hajihashemi2024multi} introduces a multi-model ensemble CP framework, optimizing model selection within a pool of experts for combined coverage.  CP has also been extended for change-point detection in time series via conformal martingales \citep{vovk2021retrain, volkhonskiy2017inductive}. 

Recent works on online conformal prediction develop online strategies that guarantee marginal coverage for adversarial or nonstationary data streams. ACI \citep{gibbs2021adaptive}, AgACI \citep{zaffran2022adaptive}, DtACI \citep{gibbs2024conformal}, and Conformal PID control \citep{angelopoulos2024pid}  leverage online optimization and have a "reaction" period when distribution shifts. Multivalid Prediction (MVP) \citep{bastani2022practical} provides group coverage guarantees and learns a calibration threshold online, although it can require longer calibration windows. Closer to our work, SPCI \citep{xu2023sequential} learns an autocorrelative estimation of the residual's quantiles through Quantile Random Forest; whereas HopCPT \citep{auer2023hopcpt} condition their conformal interval on similar parts of the time series using a Modern Hopfield Network. The drawback of these two methods is that the regression models are trained during inference time, which (1) requires a long cold start window to learn meaningful associations and (2) is computationally expensive. 
%
Our contribution is a CP algorithm that is computationally light during inference time, and can leverage state prediction capabilities within the SDS model to anticipate and adapt quickly to distribution changes.


%Closest to our work are (1) SPCI \citep{xu2023sequential} learns a time-adaptive estimation of residual quantiles over time, upon leveraging the temporal dependency among them. However, SPCI's theoretical results rely on the properties of a specific quantile learning algorithm, which does not always apply to the time series of interest.(2) Multivalid Prediction (MVP) \citep{bastani2022practical} which guarantees group conditioned coverage in the online adversarial setting. MVP requires a longer sequence as the calibration threshold is learned online. Our work brings their theoretical results into a more practical setting. Our approach also achieves group conditional coverage, which builds upon results from SPCI and MVP. A similar result can be found in \cite{zhou2024heterogeneous} with the name approximate conditional coverage, although only empirical evidence is shown.

%Most relevant to our work is the recent endeavor in generalizing CP to time-series forecasting. \cite{Schaar2021conformaltime} and \cite{sun2023copula} studies . In the case of distribution shifts, ACI \citep{gibbs2021adaptive} and EnbPI \citep{xu2021conformal} developed CP algorithms that relax the exchangeability assumption while maintaining asymptotic validity via online learning (former) and ensembling (later); \cite{zaffran2022adaptive} further improves online adaptation.  \cite{sousa2022general} combines EnbPI with conformal quantile regression \citep{romano2019conformalized} to model heteroscedastic time series. 






